\chapter{Hecke pairs and double cosets}\label{chap:hecke-pairs}

The basic combinatorial datum of Hecke theory is a \emph{Hecke pair} $(G,H,\Delta)$, from
which one builds two free $\mathbb{Z}$-modules: one on the double cosets
$H\backslash\Delta/H$, which will carry the Hecke ring structure, and one on the left
cosets $Hg$, on which that ring will act. This chapter records these definitions together
with the structural fact, following [Sh, \S3.1], that each double coset $HgH$ is a finite
disjoint union of left cosets.

\begin{definition}[Hecke pair]\label{def:hecke-pair}
\lean{HeckeRing.HeckePair}
\leanok
A \emph{Hecke pair} $(G,H,\Delta)$ consists of a group $G$, a subgroup $H$ of $G$, and a
submonoid $\Delta$ of $G$ satisfying $H\le\Delta$ and $\Delta\le\operatorname{Comm}(H)$,
where $\operatorname{Comm}(H)$ is the commensurator of $H$ in $G$. Equivalently, every
$g\in\Delta$ is commensurable with the identity in the sense that
$H\cap gHg^{-1}$ has finite index in both $H$ and $gHg^{-1}$.
\end{definition}

\begin{definition}[Double cosets]\label{def:hecke-coset}
\lean{HeckeRing.HeckeCoset}
\uses{def:hecke-pair}
\leanok
For a Hecke pair $(G,H,\Delta)$, the set of \emph{double cosets} $H\backslash\Delta/H$ is
the quotient of $\Delta$ by the equivalence relation that identifies $g$ and $g'$ whenever
$HgH=Hg'H$, that is, whenever $g'=h_1 g h_2$ for some $h_1,h_2\in H$. We write a typical
double coset as $HgH$.
\end{definition}

\begin{definition}[Left cosets in $\Delta$]\label{def:hecke-left-coset}
\lean{HeckeRing.HeckeLeftCoset}
\uses{def:hecke-pair}
\leanok
For a Hecke pair $(G,H,\Delta)$, the set of \emph{left cosets} is the quotient of $\Delta$
by the equivalence relation that identifies $g$ and $g'$ whenever $Hg=Hg'$. We write a
typical left coset as $Hg$.
\end{definition}

\begin{lemma}[Double coset as a union of left cosets]\label{lem:dc-eq-iUnion-lc}
\lean{HeckeRing.DoubleCoset.doubleCoset_eq_iUnion_leftCosets}
\uses{def:hecke-pair}
\leanok
Let $g\in\Delta$. Then the double coset $HgH$ is the union
\[
  HgH \;=\; \bigcup_{i\,\in\, H/(H\cap gHg^{-1})} H\,(g_i\,g),
\]
where $i$ ranges over the coset space $H/(H\cap gHg^{-1})$ and $g_i\in H$ is a
representative of $i$. In particular $HgH$ is a disjoint union of finitely many left
cosets of $H$.
\end{lemma}

\begin{proof}
\uses{def:hecke-pair}
\leanok
Writing $K=H\cap gHg^{-1}$, the subgroup $H$ is the disjoint union of the left cosets
$g_i K$ as $i$ ranges over $H/K$. Multiplying this decomposition on the right by the set
$gHg^{-1}$ and using that $K$ stabilises $gHg^{-1}$ under left multiplication, one obtains
$H\,gHg^{-1}=\bigcup_i g_i\,gHg^{-1}$; cancelling $g^{-1}$ on the right and translating by
$g$ turns the left-hand side into $HgH$ and each summand into the left coset $H(g_i g)$.
Commensurability of $g$ with the identity, which holds because
$\Delta\le\operatorname{Comm}(H)$, makes the index $[H:H\cap gHg^{-1}]$ finite, so the
union is finite, and distinct cosets $g_iK$ give distinct left cosets $H(g_i g)$, whence
the union is disjoint.
\end{proof}

\begin{definition}[Left-coset index of a double coset]\label{def:decomp-quot}
\lean{HeckeRing.decompQuot}
\uses{def:hecke-pair, lem:dc-eq-iUnion-lc}
\leanok
For $g\in\Delta$, the \emph{decomposition index set} of $g$ is the coset space
$H/(H\cap gHg^{-1})$, which indexes the left cosets appearing in the decomposition
$HgH=\bigsqcup_i H(g_i g)$ of Lemma~\ref{lem:dc-eq-iUnion-lc}. Since
$\Delta\le\operatorname{Comm}(H)$, this index set is finite.
\end{definition}

\begin{definition}[Hecke ring, underlying module]\label{def:hecke-ring-carrier}
\lean{HeckeRing.𝕋}
\uses{def:hecke-coset}
\leanok
Let $(G,H,\Delta)$ be a Hecke pair. The \emph{Hecke ring}
$\mathbb{T}=\mathcal{H}(G,H,\Delta)$ is, as a $\mathbb{Z}$-module, the free
$\mathbb{Z}$-module on the set of double cosets $H\backslash\Delta/H$; that is, the
$\mathbb{Z}$-valued functions on $H\backslash\Delta/H$ of finite support. Its
multiplication is introduced in Chapter~\ref{chap:ring-structure}.
\end{definition}

\begin{definition}[Hecke module]\label{def:hecke-module}
\lean{HeckeRing.HeckeModule}
\uses{def:hecke-left-coset}
\leanok
Let $(G,H,\Delta)$ be a Hecke pair. The \emph{Hecke module} is the free
$\mathbb{Z}$-module on the set of left cosets $Hg$; that is, the $\mathbb{Z}$-valued
functions of finite support on the set of left cosets. This is the module on which the
Hecke ring $\mathbb{T}$ acts.
\end{definition}
